\documentclass[12pt,a4paper]{article}

% ===== 宏包 =====
\usepackage[UTF8]{ctex}                  % 中文支持
\usepackage{fontspec}                     % 字体设置
\setmainfont{Times New Roman}             % 西文主字体
\usepackage{geometry}                     % 页面边距
\usepackage{booktabs}                     % 三线表
\usepackage{tabularx}                     % 自适应表格
\usepackage{enumitem}                     % 列表自定义
\usepackage{hyperref}                     % 超链接
\usepackage{xcolor}                       % 颜色
\usepackage{fancyhdr}                     % 页眉页脚
\usepackage{titlesec}                     % 标题格式
\usepackage{graphicx}                     % 图片
\usepackage{multirow}                     % 合并行
\usepackage{array}                        % 表格增强

% ===== 页面设置 =====
\geometry{left=2.5cm, right=2.5cm, top=2.8cm, bottom=2.5cm}

% ===== 页眉页脚 =====
\pagestyle{fancy}
\fancyhf{}
\fancyfoot[C]{\small 第 \thepage\ 页}
\renewcommand{\footrulewidth}{0.4pt}

% ===== 标题格式 =====
\titleformat{\section}{\Large\bfseries}{第\,\thesection\,节：}{0pt}{}
\titleformat{\subsection}{\large\bfseries}{\thesubsection\quad}{0pt}{}

% ===== 自定义命令 =====
\newcommand{\field}[2]{\textbf{【#1】}：\makebox[180pt]{#2}}
\newcommand{\fieldlong}[2]{\textbf{【#1】}：\makebox[380pt]{#2}}
\newcommand{\blankitem}{\item[\textbullet]\makebox[400pt]{}}

% ===== 文档开始 =====
\begin{document}

% ===== 标题 =====
\begin{center}
    {\LARGE\bfseries 课程学习小组讨论记录}\\[1cm]
\end{center}

% ===== 基本信息 =====
\noindent
\field{讨论主题}{AI 对软件开发的影响及工具分享}\\[6pt]
\field{课程名称}{软件项目开发与实践}\\[6pt]
\field{讨论日期}{\today}\\[6pt]
\field{讨论时长}{}\\[6pt]
\field{主持人}{}\\[6pt]
\fieldlong{参与人员}{}\\[12pt]

\rule{\textwidth}{0.8pt}

% ================================================================
% 第一节：讨论背景与目标
% ================================================================
\section{讨论背景与目标}

\subsection{讨论背景}
近年来，AI Coding工具（如 Cursor、CodeBuddy、Trae 等）快速发展。课程中需结合实践，思考 AI 在真实开发中的价值与风险。

% 在下方填写本次讨论的具体背景
\vspace{0.5cm}
\noindent\makebox[\textwidth]{}\\
\noindent\makebox[\textwidth]{}\\
\noindent\makebox[\textwidth]{}

\subsection{讨论目标}
\begin{enumerate}[label=\arabic*.]
    \item 分析 AI 对软件开发各环节的影响
    \item 分享常用 AI 开发工具与使用经验
    \item 总结对程序员职业发展的启示
    % 根据实际讨论需求，可增加目标条目
\end{enumerate}

% ================================================================
% 第二节：主要讨论内容
% ================================================================
\section{主要讨论内容}

\subsection{AI 对软件开发的影响}

\subsubsection{积极影响}
\begin{itemize}[leftmargin=2em]
    \item 提高编码效率（代码补全、生成、重构）
    \item 降低入门门槛（新手更容易上手）
    \item 辅助文档编写、测试用例生成
    \item 减少重复性劳动（样板代码、SQL、正则等）
    % 根据实际讨论补充
\end{itemize}

\subsubsection{潜在问题与风险}
\begin{itemize}[leftmargin=2em]
    \item 生成代码质量不稳定（Bug、逻辑错误）
    \item 安全隐患（敏感信息泄露、不安全代码）
    \item 过度依赖导致基础能力下降
    \item 版权与许可证不明确
    % 根据实际讨论补充
\end{itemize}

\subsubsection{对开发者角色的影响}
\begin{itemize}[leftmargin=2em]
    \item 从“写代码” $\rightarrow$ “审查与组合代码”
    \item 更注重系统设计、架构与业务理解
    \item 提示词工程（Prompt Engineering）成为新技能
    % 根据实际讨论补充
\end{itemize}

% 在下方补充其他讨论内容
\vspace{0.5cm}
\noindent\textbf{补充讨论内容：}\\[6pt]
\noindent\makebox[\textwidth]{}\\
\noindent\makebox[\textwidth]{}\\
\noindent\makebox[\textwidth]{}

\subsection{AI 开发工具分享}

\begin{table}[htbp]
\centering
\caption{AI 开发工具对比}
\label{tab:tools}
\renewcommand{\arraystretch}{1.3}
\begin{tabularx}{\textwidth}{|>{\centering\arraybackslash}m{2.2cm}|>{\centering\arraybackslash}m{1.8cm}|>{\centering\arraybackslash}m{2.8cm}|>{\centering\arraybackslash}m{3cm}|>{\centering\arraybackslash}m{3cm}|}
\hline
\textbf{工具名称} & \textbf{类型} & \textbf{适用场景} & \textbf{优点} & \textbf{缺点} \\
\hline
GitHub Copilot & 代码补全 & 日常编码 & 智能、流畅 & 付费、隐私 \\
\hline
ChatGPT & 通用 AI & 设计、调试、解释代码 & 灵活 & 可能产生幻觉 \\
\hline
Claude & 长文本分析 & 读源码、文档 & 上下文窗口长 & 速度一般 \\
\hline
Cursor & AI IDE & 全流程开发 & 体验好 & 学习成本 \\
\hline
% 根据实际分享补充更多工具
% 工具名 & 类型 & 场景 & 优点 & 缺点 \\
\hline
\end{tabularx}
\end{table}

% ================================================================
% 第三节：讨论总结与反思（可选）
% ================================================================
\section{讨论总结与反思}

\subsection{关键收获}
% 填写本次讨论的关键收获
\vspace{0.5cm}
\noindent\makebox[\textwidth]{}\\
\noindent\makebox[\textwidth]{}\\
\noindent\makebox[\textwidth]{}

\subsection{个人感悟}
% 填写个人感悟
\vspace{0.5cm}
\noindent\makebox[\textwidth]{}\\
\noindent\makebox[\textwidth]{}\\
\noindent\makebox[\textwidth]{}

\subsection{后续行动计划}
% 填写后续计划
\vspace{0.5cm}
\noindent\makebox[\textwidth]{}\\
\noindent\makebox[\textwidth]{}\\
\noindent\makebox[\textwidth]{}

\end{document}
